En un ensayo simulado de vaciado se registra la altura normalizada \(x(t)=H(t)/H(0)\), adimensional, con \(t\) en minutos. Se comparan dos leyes con coeficientes constantes positivos:
\[
\text{A:}\quad x'=-a x,
\qquad
\text{B:}\quad x'=-b\sqrt{x}.
\]
Ambas se aplican mientras quede líquido. La información de calibración, tratada como exacta en este ejercicio, es \(x(0)=1\) y \(x'(0)=-0.20\,\mathrm{min}^{-1}\). Se reserva para comprobación una lectura en \(t=2\) min: \(x_{\rm obs}=0.642\), con error absoluto garantizado no mayor que \(0.004\).

Determine el coeficiente de cada ley a partir de la misma calibración y obtenga sus predicciones en el instante reservado. Compruebe que se aplican antes de un posible vaciado. Decida cuáles son compatibles con la lectura y su cota de error, y explique qué permite concluir esta comprobación sobre las dos leyes y qué no establece para otros ensayos o tiempos.
